\section{Schwierigkeiten}\label{sec:probleme}

Während der Umsetzung des Konzeptes sind einige unerwartete oder nennenswerte Schwierigkeiten aufgekommen.
Diese werden nachfolgend aufgelistet (die Reihenfolge hat keine Bedeutung):

%% also good python colors: autumn, manni, perldoc, tango, lovelace, rainbow_dash, gruvbox-light
\begin{enumerateWithCode}{S}
    \item Es war nicht möglich für ein Merkmal mehrere Charakter-Fakten mit mehreren Argumenten ohne Hilfsfakten zu erzeugen.
    Zwar können \textquote{external functions} auch Listen \bzw \Tech{Clingo}-Tuple zurückgeben; diese werden aber nicht automatisch als Argumente entpackt.
    \todo{sollten das Listings sein?}Für \zB Talente sieht daher die Generierung der Merkmalsfakten wie folgt aus (vgl. \Anhang[KP]{sec:common.lp}, Zeile 63-75):
    \begin{mintedNameless}[linenos=false]{clingo}
__talent_wrapper(@talents).
talent(T,LVL) :- __talent_wrapper((T,LVL)).
    \end{mintedNameless}

    Die externe Funktion \Code{@talents} ist wie folgt umgesetzt (vgl. \Anhang[KP]{sec:hero_wrapper.py}):
    \begin{mintedNameless}{python}
% [ ... ]
def _map_skills(skills: List[Skill]) -> List[Symbol]:
    return [Tuple_([String(skill.name), Number(skill.level)]) for skill in skills]
% [ ... ]
class HeroWrapper:
    _hero: Hero
% [ ... ]
    def talents(self) -> List[Symbol]:
        return _map_skills(self._hero.talents)
% [ ... ]
    \end{mintedNameless}

    {\color{red}\item Das Definieren von sinnvollen und ausreichenden Testfällen war und ist schwierig bei den Komponententests.
    Bei den Funktionstets können einfach valide Charaktere (Testdaten) ermittelt werden und \ggf gezielt fehlerhafte Abweichungen platziert werden.
    Für die Komponententests geht dies nur bedingt, so ist es, für die Wartbarkeit, zu vermeiden die Testdaten mit unnötigen Informationen zu füllen, also Charakter-Merkmale anzugeben die für den Testfall nicht relevant sind.
    Dies erfordert wissen über die Implementierung und gegenseitige Beieinflussung von Regeln.
    Trotz der Stufen-weiße Helden-Validierung können solch gezielte Testdaten, für den Testfall, unrelevante Meldungen erzeugen.

    Es ist auch schwierig zuvermeiden ungewollt gleiches tu testen ...... Ferner, für die Wartbarkeit, ist es wichtig ein möglichst kleines Set and Testfällen zu definieren, welches dennoch eine vollständige Abdeckung der Funktionalitäten bietet und konkrete nur einzelne Funktionen prüft.

    }

    {\color{red}\item more to come}
\end{enumerateWithCode}
